%%%%%%%%%%%%%%%%%%%%%%%%%%%%%%%%%%%%%%%%%%%%%%%%%%%%%%%%%%%%%%%%%%%%%%%%%%%%%%%%
%% 
\cleardoubleoddpage%  Make sure to start each chapter on a new odd page
\chapter{Practical Analysis}
\label{sec:analysis}


\section{Ghidra}
\label{sec:analysis:ghidra}

\subsection{Installation and Setup}
\label{subsec:analysis:ghidra:setup}
Ghidra version 12.1.2 was obtained from the project's official GitHub releases page on 10 August 2026. The software was distributed as a ZIP archive 
named \texttt{ghidra\_12.1.2\_PUBLIC\_20260605.zip}. No conventional installation procedure was required. After extracting the archive, Ghidra was started 
using the included \texttt{ghidraRun} launcher. On the first launch, the Ghidra User Agreement was displayed and accepted.

A compatible Java runtime is required to run Ghidra. Java 21 is recommended for this version, although Java 25 was already installed in the evaluation 
environment and operated without problems. No administrative privileges or additional configuration were required. The application launched successfully 
immediately after extraction, and the complete setup process took only a few minutes. Apart from the requirement for an existing Java runtime, the installation 
process was straightforward.

To verify the setup, a new Ghidra project was created and a harmless synthetic sample was imported by dragging the file into the project window. 
The sample was then opened and analysed using the default analysis configuration, with all initially enabled analysis options retained. Import and automatic 
analysis completed without errors. Overall, Ghidra provided a simple setup procedure consisting only of downloading and extracting the archive before launching 
the application.

\subsection{Analysis of the Synthetic C Sample}
\label{subsec:analysis:ghidra:synthetic}

\subsection{Analysis of the WannaCry Sample}
\label{subsec:analysis:ghidra:wannacry}


\section{Binary Ninja}
\label{sec:analysis:binary-ninja}

\subsection{Installation and Setup}
\label{subsec:analysis:binary-ninja:setup}
Binary Ninja Free version 5.3.9757-Stable was obtained from the official Binary Ninja website on 10 August 2026. This is the free edition of Binary Ninja, 
rather than one of the commercially licensed versions. The Linux version was distributed as a ZIP archive named \path{binaryninja_free_linux.zip}. 
As with Ghidra, no conventional installation procedure was necessary. After downloading and extracting the archive, the application was started using the 
included \texttt{binaryninja} executable.

No additional dependencies, administrative privileges, or configuration steps were required. Binary Ninja launched successfully on the first attempt. 
During the initial launch, the application asked whether crash reports should be sent to the developers. It subsequently displayed its interface using a dark 
theme by default. The complete setup process took only a few minutes, and no problems or confusing steps were encountered.

The setup was verified by opening a harmless synthetic binary through the application's file-selection dialogue. The file was imported successfully without 
requiring additional configuration. Overall, Binary Ninja Free provided a straightforward setup procedure consisting only of downloading and extracting the 
archive before launching the included executable.

\subsection{Analysis of the Synthetic C Sample}
\label{subsec:analysis:binary-ninja:synthetic}

\subsection{Analysis of the WannaCry Sample}
\label{subsec:analysis:binary-ninja:wannacry}


\section{Cutter}
\label{sec:analysis:cutter}

\subsection{Installation and Setup}
\label{subsec:analysis:cutter:setup}
Cutter version 2.5.0 was obtained through the project's official website on 10 August 2026. The Linux version was distributed as an AppImage named 
\path{Cutter-v2.5.0-Linux-x86_64.AppImage}. After downloading the file, it was marked as executable and launched directly. No conventional installation procedure 
was required.

No additional dependencies, administrative privileges, or configuration steps were necessary. Cutter launched successfully on the first attempt. The initial 
dialogue offered the option to use the operating system's theme, although the application subsequently opened with a light theme. Apart from this minor 
inconsistency, no problems or confusing steps were encountered. The download itself accounted for most of the time required to make the application operational.

The setup was verified by selecting and opening a harmless synthetic binary through the file-selection dialogue. Cutter recognized the file and presented a slider 
for selecting the desired analysis level before starting the analysis. This provided a clear and convenient way to control the extent of the initial 
automatic analysis.

Overall, the AppImage format made the setup process particularly straightforward because the application could be launched as a self-contained executable 
without installing additional packages. Although an AppImage may introduce a small amount of overhead during application startup or when initially 
accessing compressed files, it does not involve emulation or virtualization. Its runtime performance is therefore generally comparable to that of a 
conventionally installed native application.

\subsection{Analysis of the Synthetic C Sample}
\label{subsec:analysis:cutter:synthetic}

\subsection{Analysis of the WannaCry Sample}
\label{subsec:analysis:cutter:wannacry}

\section{Bytecode Viewer}
\label{sec:analysis:bytecode-viewer}

\subsection{Installation and Setup}
\label{subsec:analysis:bytecode-viewer:setup}
Bytecode Viewer version 2.13.2 was obtained from the project's official GitHub releases page on 10 August 2026. The application was distributed as an executable JAR 
file named \path{Bytecode-Viewer-2.13.2.jar}. No conventional installation procedure was required. After downloading the file, the application was launched using 
the installed Java runtime:

\begin{lstlisting}[language=bash]
java -jar Bytecode-Viewer-2.13.2.jar
\end{lstlisting}

Apart from a compatible Java runtime, no additional dependencies were required. The application did not require administrative privileges or further configuration 
and launched successfully on the first attempt. Consequently, the complete setup required little time or effort.

The setup was verified by dragging the synthetic Java sample into the designated drop area of the application. Bytecode Viewer imported the JAR file successfully 
and decompiled its contents within a short period of time.

Bytecode Viewer was selected as a replacement for JD-GUI, which had initially been intended for the evaluation. During preliminary testing, JD-GUI repeatedly 
produced internal errors when processing the valid synthetic sample. Compiling the sample for older Java target versions, including Java 17, Java 11, and Java 8, 
did not resolve the problem. Bytecode Viewer successfully processed the sample without requiring it to be modified and provided a more actively maintained 
alternative for the practical evaluation. Overall, its installation and initial setup were straightforward and presented no difficulties.

\subsection{Analysis of the Synthetic Java Sample}
\label{subsec:analysis:bytecode-viewer:synthetic}

\section{ILSpy}
\label{sec:analysis:ilspy}

\subsection{Installation and Setup}
\label{subsec:analysis:ilspy:setup}
ILSpy version 11.0.0 was obtained from the project's official GitHub repository on 10 August 2026. The evaluated build was the release candidate distributed as the Debian package 
\path{ilspy_11.0.0.9335.rc-1_amd64.deb}. The package was installed using APT:

\begin{lstlisting}[language=bash]
sudo apt install ./ilspy_11.0.0.9335.rc-1_amd64.deb
\end{lstlisting}

Unlike the portable archives and executable JAR files used by several other evaluated tools, ILSpy required a conventional system-wide installation and therefore administrative 
privileges. The package targeted the AMD64 architecture and declared dependencies on several graphical and system libraries, including X11, ICE, SM, Fontconfig, ICU, and xdg-utils. 
APT resolved the required dependencies during installation, and no additional manual configuration was necessary.

The installation completed without errors, and ILSpy launched successfully afterward. The Debian package simplified integration into the Debian-based evaluation environment and 
made the installation process straightforward. Only little time and effort were required.

The setup was verified by importing the compiled .NET version of the synthetic sample. ILSpy recognized the managed assembly and successfully reconstructed its contents. 
Initial navigation was somewhat less intuitive because the assembly tree contained numerous referenced framework assemblies, which made the relevant application class difficult to 
locate at first. After expanding the appropriate nodes and selecting the \texttt{IbanValidator} class, the reconstructed source code could be accessed successfully. This 
represented a minor learning curve in using the interface rather than an installation or decompilation problem.

\subsection{Analysis of the Synthetic C\# Sample}
\label{subsec:analysis:ilspy:synthetic}


%% 
%%%%%%%%%%%%%%%%%%%%%%%%%%%%%%%%%%%%%%%%%%%%%%%%%%%%%%%%%%%%%%%%%%%%%%%%%%%%%%%%
